\section{F. Boundary cases, the explicit-element lead, and the audit boundary}

\begin{claim}{F.1}{Proved}{The exceptional types and centers do not permit silent transfer}
Over an algebraically closed field of characteristic zero, the geometric
fundamental groups of types $(0,0),(0,1),(0,2),(1,0)$ are, respectively,
$1,1,\widehat{\mathbb Z},\widehat{\mathbb Z}^2$.
Consequently none is a center-free hyperbolic fiber to which C.a4 applies.
For type $(0,2)$ over $\mathbb Q$ the outer action is the cyclotomic
character, and is not faithful. For an elliptic curve of type $(1,0)$,
its action on the last group is the Tate-module action of that particular
curve; it is not the action on the universal modular group $H$.
\end{claim}
\begin{proof}
The complex sphere and plane are simply connected, the punctured plane
retracts onto a circle, and a complex elliptic curve is a quotient of
$\mathbb C$ by a rank-two lattice. Finite covering theory and the
characteristic-zero comparison give the four completions. The covers
$t\mapsto t^n$ of $\mathbb G_m$ identify its arithmetic action with
the action on all roots of unity. This factors through an abelian group.
Its kernel is nontrivial: the splitting field of $t^3-2$ has Galois
group $S_3$ (irreducibility, and its nonsquare discriminant $-108$,
give this elementary cubic criterion), whereas every quotient of an
abelian group is abelian. Thus a faithful cyclotomic action would
contradict this quotient. For an elliptic curve the finite coverings
are its isogenies; the inverse system of multiplication-by-$n$ kernels
is its Tate module. This identifies the representation without importing
any assertion about the universal stack. The nontrivial centers of
the latter two groups prevent the center-free reductions used earlier.

Both $(0,3)$ and $(1,1)$ satisfy $2g-2+\nu>0$ and have free rank-two
geometric fiber groups. Their moduli and mapping-class groups are
different objects: the central $\mu_2$ in the elliptic moduli stack
and the center of $H$ were explicitly retained in A.1 and C.a7--C.a16.
At finite congruence level the center can be larger still, as C.r1
computes. The proof of D.5 uses no division, includes characteristics
two and three, and identifies the non-toric fiber at three. The Legendre
gerbe is asserted only over $\mathbb Z[1/2]$; it includes characteristic
three, but not the invalid four-etale-branch-point construction in
characteristic two.

A smooth complex curve is canonically oriented, so a nonorientable
surface is not another complex type $(g,\nu)$ to which these theorems
apply. Orientation-reversing symmetries are instead represented by
$B^\pm$ and $J$ in A.2. Over $\mathbb R$, the arithmetic quotient is
$G_\mathbb R=C_2$ and for the tripod in our normalization its outer action is represented
by the same $c=(-1,1)$; a chosen real or tangential base point supplies a genuine
lift. This covers the real involution in the manuscript without
identifying the fundamental group of a real locus with that of the
complex curve. The latter identification is generally false, since
the real locus is a union of one-dimensional manifolds.
\end{proof}

\begin{claim}{F.2}{Proved}{The explicit complex-conjugation part of Lead 3}
With a real tangential base point and the standard real access paths
to $0,1$, complex conjugation is the normalized pair $(-1,1)$.
The identity and this element are distinct full profinite solutions;
the latter has order two. These statements answer the literal examples
of \S42, fol.349\textquotesingle, scan120, and \S47, fols.426--427, scans252--253.
\end{claim}
\begin{proof}
Reflection in the real axis reverses the orientation of each small
peripheral circle while fixing the chosen real access paths. Thus its
based action is $x\mapsto x^{-1},y\mapsto y^{-1}$. The normalized
commutator conjugator is $1$, by the uniqueness proved in D.1.
This action squares to the identity and acts by minus identity on
$\widehat{\mathbb Z}^2$, so is nontrivial. Its cycle and pentagon
equations hold because it is an arithmetic action; alternatively,
the pentagon is the identity when $f=1$, and the two cycle equations
follow by direct cancellation with $\nu=-1$. The self-centralizer
calculation in B.m7 gives further information about this full
profinite element, independently of a classification of all torsion.
\end{proof}

\begin{claim}{F.3}{Open}{The exact limit of the explicit-element and full-image comparison}
We have not constructed an element of the full profinite
$\GTz\setminus\GT$, nor proved that this difference is empty.
The finite non-arithmetic outer solution of B.g15 is a solution of
the specified finite two-cycle system; it is not asserted to lift
to $\GTz$. The strongest explicit lifting result here is the
weight-seven rational solution through class twelve in B.g11.
Its first uncomputed rational extension is degree thirteen, with
the group equations and the preceding coefficients retained.
The first undecided integral extension at the prime thirteen also
leaves the range of the class-less-than-prime BCH construction.
\end{claim}
% anchor:F.3
\manuscriptanchor{\S 42; fol.349 verso; scan120}

\begin{proof}[Attack and unresolved comparison]
The raw weight-four candidate fails at degree five by the exact
rank obstruction in B.g4. The special weight-seven candidate has
a nonzero pentagon defect, and solving all the retained higher
coefficient equations simultaneously gives the class-twelve jet;
the failed greedy attempt is retained in the code. Reduction modulo
thirteen gives a genuine finite group, and the actual pure-braid
action in B.g15 proves its arithmetic exclusion. To pass from this
object to the full pro-$13$ group would require a compatible solution
at every further characteristic quotient; no inverse-limit existence
theorem follows from the twelve solved degrees. Passing from one
pro-$13$ group to the full profinite group requires further simultaneous
compatibilities. B.f8 shows concretely why automorphism tests in every
nilpotent completion cannot replace those requirements.

The sentence ``no other element is known explicitly'' is a statement
about available descriptions in a literature, not a mathematical
nonexistence assertion. We do not turn the failed construction into
such a theorem. The named comparison is Lochak and Schneps (1997),
\emph{A cohomological interpretation of the Grothendieck--Teichm\"uller
group}, Proposition~4(i),(ii), whose proof uses only (II) after the
fixed-subgroup input; see B.m7 and
\citet{LochakSchneps1997dictionary}. It proves a centralizer statement
and does not by itself classify all finite-order elements. Pro-$\ell$ torsion,
finite-shadow orders, and full profinite torsion remain distinct
questions, as B.g10 explains.
\end{proof}

\subsection*{Scope of the source and computational audit}
The accompanying report records the primary sources opened, the
specific imported theorem numbers, the original scans inspected, and
the early failed live-source requests. The checked transcriptions,
their research notes, the LM3 return and review, and the BISITE return
and review were treated as supplied leads. Our calculations were
rerun from the definitions; no assertion receives its status merely
because it had that status in an earlier return.

The bibliography marks \emph{verified} for a source opened, not for
an entire proof independently certified. The exact partial audit of
the attached 2026 preprints is B.g9. The comparator file ending in
\texttt{sorry} is not a supplied proof certificate; this statement
does not assert that a separate proof module is absent elsewhere.
The live primary arXiv record confirms the attached Willwacher
version-four date of 28 October 2026. A live request to the OpenAI
repository failed; no external formalization result is imported on
the basis of an inaccessible primary file or an unreviewed secondary
summary.

The program directory contains all mathematical programs used for
the reported calculations, their text outputs, and machine-readable
certificates. Exact integer or rational arithmetic is used; finite
field and permutation operations are checked by exhaustive enumeration
where the claims say exhaustive. Linear dimensions come from rational
rank computations and independent character formulas. A test at
finitely many levels is used as a proof of those finite cases only.
The failed and incomplete jet-solving attempts are retained and
identified as such; their failure is never used as an impossibility
proof. A deterministic build and inventory audit checks that every
numbered claim has exactly one row in the results register.

The unresolved problems are not interchangeable: E.3 concerns the
additional orbifold map, C.d5 concerns simultaneous profinite
realization of adelic normalizer coordinates, D.6 concerns conjugacy
classes of sections over a specified domain, B.f7 concerns an
arithmetic finite-shadow image, and B.g11/F.3 concern an infinite
lifting problem. Their finite attacks and smallest uncomputed
instances are stated separately. We make no claim of priority.
